
The field of large language model pre-training has developed a range of data curation methods, including perplexity-based quality filtering [Zhou et al., 2024; Gururangan et al., 2024], near-duplicate removal \citep{He2024SoftDedup}, and domain mixing [Wettig et al., 2025; Peng et al., 2025]. These methods are typically validated at a single model scale and then applied as universal recommendations regardless of model size.

A critical assumption underlying this practice is that the optimal curation configuration is independent of model capacity. Prior work has not systematically tested this assumption. Individual methods such as ProX \citep{Zhou2024ProX}, SoftDedup \citep{He2024SoftDedup}, REWIRE \citep{Nguyen2025REWIRE}, and WebOrganizer \citep{Wettig2025Organize} report improvements over baselines at specific scales, but do not vary model size as an independent factor in a controlled factorial design. The question of whether a curation recipe that benefits a 70M-parameter model also benefits a 7B-parameter model --- or whether optimal curation configurations shift with scale --- remains unaddressed in the existing literature.

This paper reports a controlled proof-of-concept experiment designed to test whether a Scale $\times$ Curation interaction exists in pre-training benchmark performance. Specifically, we ask: does the optimal GPT-2-scored perplexity filtering threshold $\tau$* depend on model size?

We find a directional signal consistent with scale-dependent optimal curation: the 7.9M-parameter model (labeled 14M) achieves peak HellaSwag performance at $\tau$\textit{=20 (strictest filtering), while the 18.1M-parameter model (labeled 31M) achieves peak performance at $\tau$}=50 (most permissive filtering). However, this result was obtained at substantially reduced scale from the original experimental target, the effect size is small, and formal statistical tests for interaction were not conducted due to insufficient power at the PoC scale.

Three main contributions are reported:

\begin{enumerate}
\item \textbf{Directional evidence of scale-dependent optimal curation}: In a 24-run factorial experiment at PoC scale ($\leq 18.1M$ parameters, 1B tokens, FineWeb), the optimal perplexity filtering threshold diverges by model size: $\tau$\textit{(7.9M) = 20, $\tau$}(18.1M) = 50.
\end{enumerate}

\begin{enumerate}
\item \textbf{A feasibility boundary for detecting this interaction}: Models at 7M/16M parameters (labeled as 7M/16M proxy in the preliminary h-e1 experiment) with 200 training steps produced no detectable signal (ANOVA p=1.0, $\eta^{2}\approx 0)$. The 14M/31M scale with 500 steps yielded the directional result.
\end{enumerate}

\begin{enumerate}
\item \textbf{Documentation of metric limitations}: MMLU 4-shot accuracy for all models below approximately 100M parameters at 200--500 training steps is at the random baseline (0.25 for 4-choice questions). HellaSwag 0-shot normalized accuracy is a more sensitive metric at this scale.
\end{enumerate}
